\chapter{Appendix}

\section{Token Verification in vm-agent}

\begin{lstlisting}[caption={[Authentication middleware]{Authentication middleware --- \filepath{vm-agent/Security/AgentTokenMiddleware.cs}}},label={lst:vm-agent-token}]
public static void UseAgentTokenAuth(this WebApplication app, string? agentToken)
{
    app.Use(async (context, next) =>
    {
        if (!string.IsNullOrWhiteSpace(agentToken))
        {
            if (!context.Request.Headers.TryGetValue("X-Agent-Token", out var provided) ||
                !TokenComparer.FixedTimeEquals(provided.ToString(), agentToken))
            {
                context.Response.StatusCode = StatusCodes.Status401Unauthorized;
                await context.Response.WriteAsJsonAsync(new
                    { detail = "Invalid or missing token (header X-Agent-Token)." });
                return;
            }
        }
        await next();
    });
}
\end{lstlisting}

\section{RunGate --- Serialised Execution}

The \texttt{RunGate} uses a \texttt{SemaphoreSlim(1,1)} with non-blocking acquisition
(\texttt{WaitAsync(0)}): if a run is already in progress, \texttt{/api/run} returns HTTP~409
without queueing a second sample.

\begin{lstlisting}[style=csharp,caption={[Single-run gate]{Single-run gate --- \filepath{vm-agent/State/RunGate.cs}}},label={lst:run-gate}]
public sealed class RunGate
{
    private readonly SemaphoreSlim _gate = new(1, 1);

    public async Task<bool> TryEnterAsync() => await _gate.WaitAsync(0);

    public void Release() => _gate.Release();
}
\end{lstlisting}

\section{Hyper-V Driver --- Agent HTTP Client}

Listing~\ref{lst:hyperv-agent} shows how the host uploads the sample and polls the behavioural
report after the VM is ready. Every request includes \texttt{X-Agent-Token} when configured.

\begin{lstlisting}[style=python,caption={[VM agent HTTP calls]{VM agent HTTP calls --- \filepath{vm\_drivers/hyperv.py}}},label={lst:hyperv-agent}]
def _agent_upload(self, job: "AnalysisJob") -> None:
    with open(job.sample_path, "rb") as f:
        files = {"file": (job.sample_path.name, f, "application/octet-stream")}
        r = requests.post(self._agent_url("/api/upload"), files=files,
                          headers=_agent_headers(), timeout=...)
    if not r.ok:
        raise RuntimeError(f"VM agent upload failed ({r.status_code})")

def _agent_get_report(self) -> Any:
    r = requests.get(self._agent_url("/api/report"), headers=_agent_headers(), ...)
    return r.json()
\end{lstlisting}

\section{VM Agent Limits}

Table~\ref{tab:vm-agent-limits} summarises the execution controls enforced by the in-guest
vm-agent through \texttt{AgentLimits.cs} and \texttt{RunGate.cs}.

\begin{table}[htbp]
    \centering
    \caption{Execution controls in vm-agent (\texttt{AgentLimits.cs}, \texttt{RunGate.cs})}
    \label{tab:vm-agent-limits}
    \begin{tabularx}{\linewidth}{@{}lY@{}}
        \toprule
        \textbf{Control} & \textbf{Value / behaviour} \\
        \midrule
        Maximum upload & 200~MB \\
        Execution timeout & Configurable; maximum 600~s \\
        Path traversal & Destination restricted to \texttt{samples/} \\
        Concurrency & One active run (\texttt{RunGate}) \\
        Authentication & \texttt{X-Agent-Token} on all routes \\
        \bottomrule
    \end{tabularx}
\end{table}

\section{Report Transfer Protocol (Path B)}

After dynamic analysis on Path~B, the guest writes the behavioural report and a completion
digest. The host copies both artefacts and accepts the report only when SHA-256 and byte length
match the guest digest. Specifically, the guest writes \filepath{C:/analysis.txt} and
\filepath{guest\_analysis\_done.txt} (JSON with \texttt{reportSha256},
\texttt{reportBytes}, \texttt{transport=psdirect}). The host copies via PsDirect or
\filepath{Copy-VMFile} and rejects truncated or tampered transfers.
